\documentclass[10pt,a4paper]{article}
\usepackage[margin=1.9cm]{geometry}
\usepackage{kotex}
\setmainhangulfont{HCR Batang}
\setsanshangulfont{HCR Dotum}
\usepackage{amsmath}
\usepackage{booktabs,tabularx,array}
\usepackage{xcolor,titlesec,enumitem}
\usepackage[hidelinks]{hyperref}
\usepackage{fancyhdr}

\definecolor{acc}{HTML}{1F5FA8}
\definecolor{boxbg}{HTML}{EEF3FA}
\titleformat{\section}{\sffamily\large\bfseries\color{acc}}{\thesection}{0.6em}{}
\titleformat{\subsection}{\sffamily\normalsize\bfseries}{\thesubsection}{0.5em}{}
\titlespacing*{\section}{0pt}{1.2em}{0.5em}
\linespread{1.2}
\setlength{\parindent}{0pt}
\setlength{\parskip}{0.45em}
\setlist{nosep,leftmargin=1.4em}
\renewcommand{\arraystretch}{1.35}
\newcolumntype{Y}{>{\raggedright\arraybackslash}X}
\newcommand{\col}[1]{\texttt{\small #1}}
\newcommand{\keybox}[1]{\par\smallskip\noindent\colorbox{boxbg}{\parbox{\dimexpr\linewidth-2\fboxsep}{\small #1}}\par\smallskip}

\pagestyle{fancy}\fancyhf{}
\fancyhead[L]{\small\sffamily 은하 SFR 예측 · 팀 실험 계획 (개정)}
\fancyhead[R]{\small\sffamily 2026-10-31}
\fancyfoot[C]{\small\thepage}

\begin{document}

{\sffamily\LARGE\bfseries\color{acc} 팀 실험 계획: 조원 자료 통합과 과적합 대책}\par
\vspace{0.3em}
{\sffamily\small 단일 모델 $\to$ 튜닝 $\to$ AutoGluon $\to$ 앙상블을 기기 4대에서 같은 기준으로 돌리는 계획}

\keybox{\textbf{요약.} (1) 모든 기기가 대회 양식의 고정 분할과 같은 5-fold, 같은 변수 코드를 쓴다. (2) 변수는 미리 정한 세트(F5, F5C, F5LS, F5X, FULL)와 fold 안에서 고른 세트(SEL)로 비교한다. (3) 변수를 늘리면 지금까지 \textbf{매번 CV가 나빠졌으므로}, 최종 선택은 ``CV 최저 모델과 1 SE 안이면 변수가 적은 모델''로 한다. (4) 검증셋은 마지막에 한 번만 본다.}

%=============================================================
\section{조원 자료를 확인한 결과}

\begin{tabularx}{\linewidth}{@{}l Y Y@{}}
\toprule
자료 & 확인한 것 & 이번 계획에서의 처리\\
\midrule
조원 1 (압축 파일) & AutoGluon 확장(14개 변수) 테스트 RMSE 8.72, 주변 은하 변수 27개, 4성분 분해·SVM 탐색. \textbf{분할이 대회 양식과 다르다}: 그 테스트 20\% 중 공식 검증셋과 겹치는 행은 20\%(8,353개)뿐 & 점수·예측값은 공식 비교와 앙상블에 쓰지 않는다. 주변 은하 27개 변수는 이미 저장소에 있다(S 묶음, 주기 경계·전체 카탈로그로 통일). 분해(hurdle) 아이디어는 AutoGluon·스태킹으로 대체\\
조원 2 (인계 문서·피처사전) & 변수 424개(25개 묶음), train 5-fold 검증에서 core(43개)가 RMSE 8.08로 가장 좋고 전부(234–328개)는 8.14–8.16으로 나빠짐 & \col{sfr\_features.py}를 받아 \col{src/}에 넣었다. 원본 13열을 뺀 411개가 \col{t2\_} 접두사로 추가된다. T2CORE = 원본 13 + core 28 = 41개, FULL = F5X 153 + 411 = 564개. 조원 2의 SHAP 선별 89개는 학습 행 절반에서 골라 CV가 낙관적이므로 제외하고, fold 안 선택 SEL(FULL)로 대신한다\\
\bottomrule
\end{tabularx}

두 조원 모두 변수를 늘리면 나빠지는 결과를 얻었고, 우리 실험도 같다(4절). 과적합 우려는 실제로 확인된 문제다.

%=============================================================
\section{과적합 대책}
\begin{enumerate}
\item \textbf{세트를 미리 정한다.} 결과를 보기 전에 정한 세트만 비교한다: F5(25, 원래 최종), F5C(47, 문헌 근거가 가장 강한 L1·L3·L4·L7), F5LS(89, 문헌 + 조원 1), F5X(153, 우리 변수 전부), FULL(우리 + 조원 2), T2CORE(조원 2 core).
\item \textbf{데이터로 고를 때는 fold 안에서만.} SEL 세트는 바깥 fold마다 학습부 안에서 Boruta 방식(각 변수를 뒤섞은 가짜 사본보다 SHAP 중요도가 커야 남김, 안쪽 3-fold 모두에서)과 중복 제거(순위 상관 0.98 초과)로 고른다. 평가 fold는 변수 선택에 관여하지 않는다.
\item \textbf{튜닝도 묶음 단위로.} 트리 튜닝에서 문헌 L1–L7과 조원 1 변수 S를 묶음별로 켜고 끄는 것을 Optuna가 nested로 고른다.
\item \textbf{규제.} 큰 세트에는 규제 LightGBM(잎 15, 최소 표본 100, 열 샘플링 0.5, L2 5)을 함께 돌린다.
\item \textbf{최종 선택 규칙(1-SE).} CV 최저 모델과 짝지은 차이가 1 SE 안인 모델 중 \textbf{입력 변수가 가장 적은} 것을 고른다. 앙상블은 가장 복잡한 것으로 친다.
\item \textbf{검증셋 1회.} 공식 검증셋은 A가 최종 선택 뒤에 한 번만 평가한다.
\end{enumerate}

%=============================================================
\section{모델 실행 체계}
\begin{tabularx}{\linewidth}{@{}l Y l@{}}
\toprule
단계 & 내용 & 세트\\
\midrule
단일 모델 & LightGBM(기본·규제), XGBoost, CatBoost(Tweedie), RF, ExtraTrees, MLP, 로그 목표 LightGBM & F5, F5C, F5LS\\
선형·커널·KNN & Ridge, Lasso, Tweedie GLM(스플라인), 커널 Ridge, KNN — 튜닝 없음 / nested 튜닝 & F5, F5C, F5LS\\
딥러닝·사전학습 & Tweedie MLP, TabM, RealMLP(원·로그 목표), xRFM, TabPFN v2; TabM·RealMLP nested 튜닝 & F5, F5C, F5LS / 튜닝은 F5C\\
트리 튜닝 & LightGBM·XGBoost·CatBoost nested Optuna: 모델 + 물리적 선택(Tweedie $p$, 반경, 문턱, 거대 은하 가중치) + 묶음 포함 여부 & F5 + 선택 묶음\\
변수 선택 & fold 안 Boruta 선택 + LightGBM·규제·XGBoost·CatBoost & SEL(F5X), SEL(FULL)\\
AutoGluon & 변수 축: F5(완료), F5C, F5LS, F5X, FULL, T2CORE. 전처리 축(F5C): Tweedie 추가, 로그·제곱근 목표, 거대 은하 가중치, extreme\_quality & 표 참조\\
앙상블 & 위 모든 OOF로 가중 평균, greedy, 상위 $k$, 스태킹(음 아닌 선형, 작은 LightGBM). 모두 CV 안에서 학습 & 전부\\
\bottomrule
\end{tabularx}

%=============================================================
\section{지금까지 확인한 결과 (CV RMSE, 원래 최종 R0 = 7.401)}
\begin{tabularx}{\linewidth}{@{}l r Y@{}}
\toprule
실험 & CV & R0 대비\\
\midrule
AutoGluon, F5 & 7.317 & $-0.08$ (SE 0.17) — 유일하게 앞섬, 노이즈 범위\\
LightGBM + L7(묶인 동반 은하) & 7.396 & 거의 같음\\
LightGBM + 다른 문헌 묶음 하나씩 & 7.51–7.62 & $+0.11$–$+0.22$\\
LightGBM + 조원 1 변수(S) & 7.562 & $+0.16$\\
LightGBM + 문헌 전부 / 문헌 + 조원 1 & 7.94 / 7.84 & $+0.54$ / $+0.44$\\
fold 안 Boruta 선택(153개 중 46–66개) & 7.65–7.84 & $+0.25$–$+0.44$ (모델 4종)\\
XGBoost, CatBoost, RF, MLP 단독(F5) & 7.62–8.60 & 모두 나쁨\\
선형·GLM(F5) & 8.70–10.0 & 크게 나쁨\\
\bottomrule
\end{tabularx}
튜닝 없는 모델에서는 변수를 늘리면 일관되게 나빠졌다. 남은 질문은 튜닝·AutoGluon·딥러닝처럼 설정이 변수에 맞춰지는 방법에서도 그런가이다.

%=============================================================
\section{기기별 작업}
\begin{tabularx}{\linewidth}{@{}l l Y r@{}}
\toprule
기기 & 작업 & 내용 & 시간\\
\midrule
A 노트북 & J01 & 단일 모델 15개(F5C 8, F5LS 7) & 2.0 h\\
 & J02 & 선형·커널·KNN 30개(3세트 $\times$ 5모델 $\times$ 튜닝 유무) & 4.5 h\\
 & J20 & fold 안 변수 선택: SEL(F5X)(완료), SEL(FULL) & 1.0 h\\
 & J21–J22 & AutoGluon F5C 로그 목표, 제곱근 목표 & 3.0 h\\
 & 마지막 & 합치기, 앙상블·스태킹, 1-SE 최종 선택, 검증셋 1회, 보고서 & 2 h\\
\midrule
B Mac & J07 & 딥러닝·사전학습 18개(6모델 $\times$ F5, F5C, F5LS) & 6.0 h\\
 & J19 & TabM·RealMLP nested 튜닝(F5C) & 5.0 h\\
\midrule
C CPU & J09–J11 & LightGBM 50회·XGBoost 40회·CatBoost 25회 nested 튜닝 & 9.0 h\\
 & J23 & AutoGluon F5C 거대 은하 가중치 & 1.5 h\\
\midrule
D GPU & J12–J15 & AutoGluon F5C, F5LS, F5X, FULL, T2CORE & 7.5 h\\
 & J16–J17 & AutoGluon F5C Tweedie 추가, extreme\_quality & 3.5 h\\
\midrule
X 선택 & J30–J31 & AutoGluon F0·F3, F5LS 4시간판 (먼저 끝난 기기) & 3–7 h\\
\bottomrule
\end{tabularx}
권장 배정: B = 조원 2(Mac), D = 조원 1(Colab GPU), C = 조원 3, A = 이 노트북.

%=============================================================
\section{실행 순서}
\begin{enumerate}
\item 모든 기기: 같은 커밋 확인, 환경 설치, CSV 배치 $\to$ \col{python scripts/jobs.py run <기기>}.
\item 각 기기: \col{python scripts/jobs.py pack <기기>} $\to$ zip을 A에 직접 전달(공개 금지).
\item A: \col{jobs.py unpack} $\to$ \col{a04\_ensemble.py} $\to$ 보고서·GitHub 갱신.
\end{enumerate}

\end{document}
